\exercisegroup

\begin{enumerate}
\def\labelenumi{\arabic{enumi}.}
\item
  设 E 是一个有序集，其中至少存在一对不同的可比较元素。证明：若 R \(\{x, y\}\) 表示关系“ \(x \in E\) 并且 \(y \in E\) 并且 \(x < y\) ''，则R满足第1号的前两个条件，但不满足第三个条件。
\item
  \begin{enumerate}
  \def\labelenumii{(\alph{enumii})}
  \tightlist
  \item
    设 E 是一个预序集，且设 S\{x, y\} 是 E 上的一个等价关系。设 R\{X, Y\} 表示关系“X ∈ E/S 且 Y ∈ E/S，且对每个 x ∈ X，存在 y ∈ Y 使得 x ≤ y”。证明 R 是 E/S 上的预序关系，称为关系 x ≤ y 关于 S 的商。集合 E/S 赋以该预序关系后，称为（按语言习惯的滥用；参见第四章第2节第6号）预序集 E 关于 S 的商。
  \end{enumerate}
\end{enumerate}

\begin{enumerate}
\def\labelenumi{(\alph{enumi})}
\setcounter{enumi}{1}
\tightlist
\item
  让 \(\varphi\) 为E到E/S上的典范映射。证明若 \(g\) 是从预序商集E/S到预序集F的一个映射，使得 \(g \circ \varphi\) 是递增映射，则 \(g\) 是递增映射。映射 \(\varphi\) 是递增的，当且仅当 S 满足以下条件：
\end{enumerate}

\begin{enumerate}
\def\labelenumi{(\Alph{enumi})}
\setcounter{enumi}{2}
\tightlist
\item
  关系 \(x \leqslant y\) 并且 \(x \equiv x' \pmod{S}\) 在 \(E\) 中蕴含存在 \(y' \in E\) 使得 \(y \equiv y' \pmod{S}\) 并且 \(x' \leqslant y'\) .
\end{enumerate}

如果满足此条件，则称等价关系 S 与预序关系 \(x \leqslant y\) （在 x 和 y 上）弱相容。每一个与预序关系 \(x \leqslant y\) （在 x 上）相容的等价关系 S （第二章，§6，第3号）更是（在x和y上）与该关系弱相容。

\begin{enumerate}
\def\labelenumi{(\alph{enumi})}
\setcounter{enumi}{2}
\item
  让 \(E_{1}\) 和 \(E_{2}\) 是两个预序集。证明如果 \(S_{1}\) 是等价关系 \(pr_{1}z = pr_{1}z'\) （在 \(E_{1} \times E_{2}\) 上），那么 \(S_{1}\) 在 \(z\) 和 \(t\) 中与乘积预序关系 \(z \leqslant t\) （\(\mathbf{E}_1 \times \mathbf{E}_2\) 上的）是弱相容的（但通常与这个关系在 \(z\) 或 \(t\) 中单独地并不相容）；此外，如果 \(\varphi_1\) 是从 \(\mathbf{E}_1 \times \mathbf{E}_2\) 到 \((\mathbf{E}_1 \times \mathbf{E}_2)/\mathbf{S}_1\) 上的典范映射，并且如果 \(\mathrm{pr}_1 = f_1 \circ \varphi_1\) 是 \(\mathrm{pr}_1\) 关于等价关系 \(\mathbf{S}_1\) 的典范分解，那么 \(f_1\) 是从 \((\mathbf{E}_1 \times \mathbf{E}_2)/\mathbf{S}_1\) 到 \(\mathbf{E}_1\) 上的一个同构。
\item
  在(a)的假设下，设E是一个有序集，并假设满足以下条件：
\end{enumerate}

(C') 关系 \(x \leqslant y \leqslant z\) 并且 \(x \equiv z \pmod{S}\) 在 E 中蕴含 \(x \equiv y \pmod{S}\) .

证明 \(\mathbf{R}\{\mathbf{X},\mathbf{Y}\}\) 则是 \(\mathbf{X}\) 与 \(\mathbf{Y}\) 之间的一个序关系（在 \(\mathbf{E} / \mathbf{S}\) 上）。

\begin{enumerate}
\def\labelenumi{(\alph{enumi})}
\setcounter{enumi}{4}
\item
  给出一个具有四个元素的全序集 E 以及 E 上的一个等价关系 S 的例子，使得条件 (C) 和 (C') 均不满足，但 E/S 是一个有序集。
\item
  设E是一个有序集，令f为从E到有序集F的递增映射，且令 \(S\{x,y\}\) 为等价关系 \(f(x)=f(y)\) 在 E 上。则条件 \((\mathbf{C}')\) 得到满足。此外，条件(C)成立当且仅当这些关系 \(x\leqslant y\) 并且 \(f(x)=f(x')\) 意味着存在 \(y'\in E\) 使得 \(x'\leqslant y'\) 并且 \(f(y)=f(y')\) . 设 \(f=g\circ\varphi\) 是 f 的典范分解。那么 g 是从 E/S 到 \(f(\mathbf{E})\) 上的一个同构，当且仅当此条件满足，并且此外，关系 \(f(x)\leqslant f(y)\) 蕴含存在 \(x',y'\) 使得 \(f(x)=f(x')\) , \(f(y)=f(y')\) ，以及 \(x'\leqslant y'\) .
\end{enumerate}

\begin{enumerate}
\def\labelenumi{\arabic{enumi}.}
\setcounter{enumi}{2}
\tightlist
\item
  令 I 为一个有序集，且令 \((\mathbf{E}_i)_{i\in \mathbf{I}}\) 是 I 指标化的一族非空有序集。
\end{enumerate}

\begin{enumerate}
\def\labelenumi{(\alph{enumi})}
\item
  设 F 为该族 \((\mathbf{E}_t)_{t \in I}\) 的和（第二章，§4，第 8 号）；对于每个 \(x \in F\) ，让 \(\lambda(x)\) 为指标 \(t\) （满足 \(x \in \mathbf{E}_t\) ）；并设 G 为由所有满足以下条件的对 \((x, y) \in F \times F\) 组成的图：要么 \(\lambda(x) < \lambda(y)\) ，要么 \(\lambda(x) = \lambda(y)\) 并且 \(x \leqslant y\) 在 \(\mathbf{E}_{\lambda(x)}\) 中。证明G是F上一个序的图。赋予此序的集合F称为该族 \((\mathbf{E}_t)_{t \in I}\) 的序数和（相对于 I 上的序）并记为 \(\sum_{t \in I} \mathbf{E}_t\) 。证明与F的划分 \((\mathbf{E}_t)_{t \in I}\) 对应的等价关系满足习题2的条件(C)和(C')，并且商有序集（习题2）典范同构于I。
\item
  如果集合I是由有序集构成的族 \((J_{\lambda})_{\lambda\in L}\) 的序数和，其中 L 是一个有序集，证明该有序集 \(\sum_{i\in I}E_{i}\) 典范同构于序数和 \(\sum_{\lambda\in L}F_{\lambda}\) ，其中 \(F_{\lambda}=\sum_{i\in I_{\lambda}}E_{i}\) （序数和的“结合律”）
\end{enumerate}

若 I 是线性有序集 \(\{1, 2\}\) ，我们记 \(E_{1} + E_{2}\) 以表示 \(E_{1}\) 和 \(E_{2}\) 的序数和。证明 \(E_{2} + E_{1}\) 和 \(E_{1} + E_{2}\) 不一定是同构的。

\begin{enumerate}
\def\labelenumi{(\alph{enumi})}
\setcounter{enumi}{2}
\item
  序数和 \(\sum_{i\in I}E_i\) 是右有向的当且仅当 I 是右有向的，且 \(E_{\omega}\) （对 I 的每个极大元素 \(\omega\) ）都是右有向的。
\item
  序数和 \(\sum_{i\in I}E_i\) 是全序的当且仅当 \(I\) 与每个 \(E_i\) 都是全序的。
\item
  序数和 \(\sum_{i\in I}E_i\) 是一个格，当且仅当满足以下条件：
\end{enumerate}

\begin{enumerate}
\def\labelenumi{(\Roman{enumi})}
\item
  集合 I 是一个格，并且对于每一对 I 中不可比较的指标 \((\lambda, \mu)\) ， \(E_{\text{sup}(\lambda, \mu)}\) （相应地， \(E_{\text{inf}(\lambda, \mu)}\) ）有一个最小（相应地，最大）元素。
\item
  对于每一个 \(\alpha \in I\) 和每一对元素 \((x, y)\) （\(E_{\alpha}\) 的元素）：若集合 \(\{x, y\}\) 在 \(E_{\alpha}\) 中有上界（相应地，有下界），则集合 \(\{x, y\}\) 在 \(E_{\alpha}\) 中有最小上界（相应地，有最大下界）。
\item
  对于每一个 \(\alpha \in I\) （\(E_{\alpha}\) 包含一个在 \(E_{\alpha}\) 中没有上界（相应地，没有下界）的两元素集合），指标 \(\lambda \in I\) （满足 \(\lambda > \alpha\) （相应地， \(\lambda < \alpha\) ））所成的集合有一个最小元素（相应地，最大元素） \(\beta\) ，以及 \(E_{\beta}\) 有最小元素（相应地，最大元素）。
\end{enumerate}

* 4. 设 E 是一个有序集，且令 \((\mathbf{E}_t)_{t \in I}\) 为E的划分，由E的连通分量（第二章，§6，习题10）关于自反且对称的关系“要么 \(x = y\) 要么 \(x\) 和 \(y\) 是不可比较的”。

\begin{enumerate}
\def\labelenumi{(\alph{enumi})}
\item
  证明如果 \(\iota \neq \kappa\) 并且 \(x \in E_{\iota}\) 并且 \(y \in E_{\kappa}\) ，那么 \(x, y\) 是可比较的；并且如果，例如， \(x \leqslant y\) , \(y' \in E_{\kappa}\) ，以及 \(y' \neq y\) ，那么也有 \(x \leqslant y'\) （利用这样一个事实：不存在将 \(E_{\kappa}\) 划分为两个集合A和B的划分，使得A中的每个元素都与B中的每个元素可比较）。
\item
  由(a)推出，与E的划分 \((\mathbf{E}_{t})\) 相对应的等价关系S在x和y上与E上的序关系 \(x \leqslant y\) 是相容的，并且商有序集E/S（习题2）是全序的。
\item
  乘积为两个全序集的有序集 \(\mathbf{E} = \mathbf{F} \times \mathbf{G}\) 的连通分支是什么？*
\end{enumerate}

\begin{enumerate}
\def\labelenumi{\arabic{enumi}.}
\setcounter{enumi}{4}
\item
  设E是一个有序集。称E的子集X是自由的，如果X中没有两个不同的元素是可比较的。设 \(\Im\) 为 E 的自由子集的集合。证明在 \(\Im\) 上，关系“对任意 \(x \in X\) , 存在 \(y \in Y\) 使得 \(x \leqslant y\) ”是X与Y之间的序关系，记作 \(X \leqslant Y\) 。映射 \(x \to \{x\}\) 是E到有序集 \(\mathfrak{J}\) 的一个子集上的同构。如果 \(X \subset Y\) ，其中 \(X \in \mathfrak{J}\) 并且 \(Y \in \mathfrak{J}\) ，证明 \(X \leqslant Y\) 。有序集 \(\mathfrak{J}\) 是全序的当且仅当 \(E\) 是全序的，并且此时 \(\mathfrak{J}\) 典范同构于 \(E\) 。
\item
  设 E 和 F 是两个有序集，并设 \(\mathcal{A}(\mathbf{E},\mathbf{F})\) 为乘积有序集 \(F^{E}\) 的由E到F的递增映射组成的子集。
\end{enumerate}

\begin{enumerate}
\def\labelenumi{(\alph{enumi})}
\item
  证明：若E、F、G是有序集，则有序集 \(\mathcal{A}(\mathbf{E}, \mathbf{F} \times \mathbf{G})\) 与乘积有序集 \(\mathcal{A}(\mathbf{E}, \mathbf{F}) \times \mathcal{A}(\mathbf{E}, \mathbf{G})\) 同构。
\item
  证明：若E、F、G是三个有序集，则有序集 \(\mathcal{A}(\mathbf{E} \times \mathbf{F}, \mathbf{G})\) 与有序集 \(\mathcal{A}(\mathbf{E}, \mathcal{A}(\mathbf{F}, \mathbf{G}))\) 同构。
\item
  如果 \(\mathbf{E} \neq \emptyset\) ，那么 \(\mathcal{A}(\mathbf{E}, \mathbf{F})\) 是格当且仅当 \(\mathbf{F}\) 是格。
\item
  假设E和F均非空。则 \(\mathcal{A}(\mathbf{E}, \mathbf{F})\) 是全序的，当且仅当下列条件之一成立：
\end{enumerate}

(α) F 由单个元素组成；

(β) E 由单个元素组成且 F 是全序的；

(γ) E 和 F 都是全序的且 F 有两个元素。

\begin{enumerate}
\def\labelenumi{\arabic{enumi}.}
\setcounter{enumi}{6}
\item
  为了使从有序集 E 到具有至少两个元素的有序集 F 的每一个既是递增映射又是递减映射的映射在 E 上恒为常数，必要且充分条件是 E 关于自反且对称的关系“x 与 y 可比较”（第二章，§ 6，习题 10）是连通的。若 E 是左有向或右有向的，则此条件特别地得到满足。
\item
  设 E 和 F 是两个有序集，设 f 是 E 到 F 的递增映射，g 是 F 到 E 的递增映射。设 A（相应地 B）是所有元素 \(x \in E\) （相应地， \(y \in F\) ）所成的集合，这些元素满足 \(g(f(x)) = x\) （相应地， \(f(g(y)) = y\) ）。证明两个有序集 A 和 B 是典范同构的。
\end{enumerate}

* 9. 若 E 是一个格，证明

\[
\sup _ {j} \left(\inf _ {i} x _ {i j}\right) \leqslant \inf _ {i} \left(\sup _ {j} x _ {i j}\right)
\]

对于每一个有限的“双重”族 \((x_{ij})\) . *

\begin{enumerate}
\def\labelenumi{\arabic{enumi}.}
\setcounter{enumi}{9}
\tightlist
\item
  设 E 和 F 是两个格。则从 E 到 F 的映射 f 是递增的，当且仅当
\end{enumerate}

\[
f (\inf (x, y)) \leqslant \inf (f (x), f (y))
\]

对于所有 \(x \in \mathbf{E}\) 以及所有 \(y \in \mathbf{E}\) .

* 给出一个递增映射 \(f\) 的例子（从乘积有序集 \(\mathbf{N} \times \mathbf{N}\) 到有序集 \(\mathbf{N}\) 中），使得关系

\[
f (\inf (x, y)) = \inf (f (x, y))
\]

对至少一对 \((x, y) \in \mathbf{N} \times \mathbf{N}_{*}\) 来说为假。

\begin{enumerate}
\def\labelenumi{\arabic{enumi}.}
\setcounter{enumi}{10}
\tightlist
\item
  格E被称为完备的，如果E的每个子集在E中都有最小上界和最大下界；这特别意味着E具有最大元素和最小元素。
\end{enumerate}

\begin{enumerate}
\def\labelenumi{(\alph{enumi})}
\item
  证明：如果一个有序集 \(\mathbf{E}\) 使得 \(\mathbf{E}\) 的每一个子集在 \(\mathbf{E}\) 中都有最小上界，那么 \(\mathbf{E}\) 是一个完全格。
\item
  有序集合的乘积是完备格当且仅当每个因子都是完备格。
\item
  序数和（习题3） \(\sum_{i\in I}E_i\) 是一个完全格，当且仅当满足以下条件：
\end{enumerate}

\begin{enumerate}
\def\labelenumi{(\Roman{enumi})}
\item
  I 是一个完全格。
\item
  如果 \(J\) 是 \(I\) 的一个没有最大元素的子集，并且如果 \(\sigma = \sup J\) ，那么 \(E_{\sigma}\) 有最小元素。
\item
  对于每一个 \(t \in I\) ， \(E_t\) 的每个在 \(E_t\) 中有上界的子集在 \(E_t\) 中都有最小上界。
\item
  对于每一个 \(\iota\in I\) （\(E_{\iota}\) 没有最大元素），由所有 \(x>\iota\) 组成的集合有一个最小元素 \(\alpha\) ，以及 \(E_{\alpha}\) 有最小元素。
\end{enumerate}

\begin{enumerate}
\def\labelenumi{(\alph{enumi})}
\setcounter{enumi}{3}
\tightlist
\item
  有序集 \(\mathcal{A}(\mathbf{E},\mathbf{F})\) （由有序集 E 到有序集 F 中的递增映射全体组成；习题 6）是完全格当且仅当 F 是完全格。
\end{enumerate}

\begin{enumerate}
\def\labelenumi{\arabic{enumi}.}
\setcounter{enumi}{11}
\item
  让 \(\Phi\) 为由集合 A 到自身的一族映射。令 \(\mathfrak{F}\) 为 \(\mathfrak{P}(A)\) 中由所有满足以下条件的 \(X \subset A\) 组成的子集：有 \(f(X) \subset X\) （对每一个 \(f \in \Phi\) ）。证明 \(\mathfrak{F}\) 关于包含关系而言是一个完全格。
\item
  设E是一个有序集。将E映射到自身的映射f被称为闭包，如果它满足以下条件：(1) f是递增的；(2) 对于每个 \(x \in E\) , \(f(x) \geqslant x\) ； (3) 对每个 \(x \in E\) , \(f(f(x)) = f(x)\) 。设 F 为 E 中在 f 下不变的元素的集合。
\end{enumerate}

\begin{enumerate}
\def\labelenumi{(\alph{enumi})}
\item
  证明对于每个 \(x \in E\) ，集合 \(F_{x}\) （由元素 \(y \in F\) 满足 \(x \leqslant y\) 者组成）非空且有最小元素，即 \(f(x)\) 。反之，若G是E的子集，使得对每个 \(x \in E\) ，由所有元素 \(y \in G\) （满足 \(x \leqslant y\) ）组成的集合有一个最小元素 \(g(x)\) ，那么g是一个闭包，而G是E中在g下不变的元素所构成的集合。
\item
  设E是一个完全格。证明F的任意非空子集在E中的最大下界属于F。
\item
  证明若 E 是一个格，则 \(f(\sup(x,y)) = f(\sup(f(x),y))\) 对于 E 中每一对元素 x, y。
\end{enumerate}

\begin{enumerate}
\def\labelenumi{\arabic{enumi}.}
\setcounter{enumi}{13}
\item
  设 A 和 B 是两个集合，并设 R 是 \(A \times B\) 的任意子集。对于 A 的每个子集 X（相应地，B 的每个子集 Y），设 \(\rho(X)\) （相应地， \(\sigma(Y)\) ）表示所有使得 \((x, y) \in R\) 对于所有 \(x \in A\) （相应地， \((x, y) \in R\) 对于所有 \(y \in B\) ）成立的 \(y \in B\) （相应地， \(x \in A\) ）所组成的集合。证明 \(\rho\) 和 \(\sigma\) 是递减映射，并且映射 \(X \to \sigma(\rho(X))\) 和 \(Y \to \rho(\sigma(Y))\) 是 \(\mathfrak{P}(A)\) 和 \(\mathfrak{P}(B)\) 中的闭包（练习13）（分别按包含关系排序）。
\item
  \begin{enumerate}
  \def\labelenumii{(\alph{enumii})}
  \tightlist
  \item
    设E是一个有序集，并且对于E的每个子集X，令 \(\rho(X)\) （相应地， \(\sigma(X)\) ）表示 X 在 E 中的上界（相应地，下界）集合。证明，在 \(\mathfrak{P}(E)\) 中，集合 \(\tilde{E}\) （由满足 \(X = \sigma(\rho(X))\) 的子集 X 组成）是一个完全格，并且映射 \(i: x \to \sigma(\{x\})\) 是 E 到有序子集 \(E'\) （\(\tilde{E}\) 的子集）上的（称为典范的）同构，使得，如果 E 的元素的一族（ \(x_t\) ）的并（相应地，交）在 E 中存在，则此并（相应地，交）的像是在 \(\tilde{E}\) 中诸 \(x_t\) 的像所成的族的并（相应地，交）。 \(\tilde{E}\) 称为有序集 E 的完备化。
  \end{enumerate}
\end{enumerate}

\begin{enumerate}
\def\labelenumi{(\alph{enumi})}
\setcounter{enumi}{1}
\item
  证明：对于E的每一个子集X， \(\sigma(\rho(X))\) 是 \(\tilde{E}\) 的子集 \(i(X)\) 在 \(\tilde{E}\) 中的上确界。若 f 是 E 到完备格 F 中的任意递增映射，则存在唯一的递增映射 \(\bar{f}\) （从 \(\tilde{E}\) 到 F 中），使得 \(f = \bar{f} \circ i\) 并且 \(\bar{f}(\sup Z) = \sup (\bar{f}(Z))\) 对于 \(\tilde{E}\) 的每一个子集 Z。
\item
  若 E 是全序的，证明 \(\tilde{E}\) 是全序的。
\end{enumerate}

¶ 16. 格 E 称为分配格，如果它满足以下两个条件：

\[
\begin{array}{l l} \left(\mathbf {D} ^ {\prime}\right) & \sup (x, \inf (y, z)) = \inf (\sup (x, y), \sup (x, z)), \\ \left(\mathbf {D} ^ {\prime \prime}\right) & \inf (x, \sup (y, z)) = \sup (\inf (x, y) \inf (x, z)) \end{array}
\]

对于所有 \(x, y, z\) 在 E 中。全序集是分配格。

\begin{enumerate}
\def\labelenumi{(\alph{enumi})}
\tightlist
\item
  证明条件 (D) 和 (D') 各自单独蕴含条件
\end{enumerate}

\[
\begin{array}{l} \text {(D)} \sup (\inf (x, y), \inf (y, z), \inf (z, x)) \\ = \inf (\sup (x, y), \sup (y, z), \sup (z, x)) \\ \text {对 所有 x,y,z 中 E.} \end{array}
\]

\begin{enumerate}
\def\labelenumi{(\alph{enumi})}
\setcounter{enumi}{1}
\tightlist
\item
  证明条件 (D) 蕴含条件
\end{enumerate}

\[
\text {   若   } x \geqslant z, \text {   则   } \sup (z, \inf (x, y)) = \inf (x, \sup (y, z)).\tag{M}
\]

由此推出(D)蕴含(D')和(D'\,`)，进而三个公理(D)、(D')和(D'\,`)等价（例如，为证明(D)蕴含(D')，取x与(D)两边的最小上界，并利用(M)）。

\begin{enumerate}
\def\labelenumi{(\alph{enumi})}
\setcounter{enumi}{2}
\tightlist
\item
  证明以下两个条件中的每一个
\end{enumerate}

\[
\begin{array}{r l} \inf (z, \sup (x, y)) & \leqslant \sup (x, \inf (y, z)), \\ \inf (\sup (x, y), \sup (z, \inf (x, y))) & = \sup (\inf (x, y), \inf (y, z), \inf (z, x)) \end{array}
\]

对于所有 \(x, y, z\) 在 \(\mathbf{E}\) 都是 \(\mathbf{E}\) 成为分配格的充分必要条件。（为证明 \((\mathbf{T}')\) 蕴含 \((\mathbf{D}^n)\) ，考虑元素

\[
\inf (z, \sup (x, \inf (y, z))).)
\]

\begin{enumerate}
\def\labelenumi{\arabic{enumi}.}
\setcounter{enumi}{16}
\tightlist
\item
  一个具有最小元素 \(\alpha\) 的格E被称为相对有补的，如果对于每一对元素 \(x, y\) （满足 \(x \leqslant y\) ），存在一个元素 \(x'\) 使得 \(\sup (x, x') = y\) 并且 \(\inf (x, x') = \alpha\) 。这样的元素 \(x'\) 称为 \(x\) 关于 \(y\) 的相对补。
\end{enumerate}

（a）证明：维数 ≥ 2 的向量空间的子空间集合 E，按包含关系排序，是一个相对有补格；但如果 x, y 是 E 中满足 x ≤ y 的两个元素，则一般而言，x 关于 y 存在多个不同的相对补。

\begin{enumerate}
\def\labelenumi{(\alph{enumi})}
\setcounter{enumi}{1}
\item
  若E是分配的且相对有补的，证明若 \(x \leqslant y\) 在E中，x关于y存在唯一的相对补。若E是分配的且相对有补的，并且此外它还有最大元 \(\omega\) ，则称E为布尔格。对于每一个 \(x \in E\) ，让 \(x^{*}\) 为 x 关于 \(\omega\) 的补。则映射 \(x \to x^{*}\) 是从 E 到通过赋予 E 相反序关系所得的序集上的一个同构，并且我们有 \((x^{*})^{*} = x\) 。若 A 是任意集合，则 A 的所有子集所成的集合 \(\mathfrak{P}(A)\) （按包含关系排序）是一个布尔格。
\item
  若E是完全布尔格（习题11），证明对 E 的元素的每个族 \((x_{\lambda})\) 以及每个 \(y \in E\) ，我们拥有
\end{enumerate}

\[
\inf (y, \sup _ {\lambda} (x _ {\lambda})) = \sup _ {\lambda} (\inf (y, x _ {\lambda})).
\]

（化简为 \(y = \alpha\) 的情形，并利用以下事实：如果 \(\inf (z, x_{\lambda}) = \alpha\) （对每一个索引 \(\lambda\) ），那么 \(z^{*} \geqslant x_{\lambda}\) （对于每一个 \(\lambda\) ）。

¶ * 18. 设 A 是一个至少含有三个元素的集合，令 ℜ 为 A 的所有划分构成的集合，按关系“ω 比 ω' 更细”\,（此关系定义在ω与ω'之间；第1号，例4）来排序。证明ℜ是一个完全格（习题11），不是分配的（习题17），但是相对有补的。（为证明最后一个断言，将属于一个划分的集合良序化。）*

\begin{enumerate}
\def\labelenumi{\arabic{enumi}.}
\setcounter{enumi}{18}
\tightlist
\item
  一个有序集 E 被称为无间隙的，如果它包含两个不同的可比较元素，并且对于每一对元素 x, y，只要 x \textless{} y，开区间 {]}x, y{[} 非空。证明一个序数和 \(\sum_{i\in I}E_{i}\) （习题 3）是无间隙的当且仅当以下条件成立：
\end{enumerate}

\begin{enumerate}
\def\labelenumi{(\Roman{enumi})}
\item
  要么 I 包含两个不同的可比较元素，要么存在 \(\iota \in I\) 使得 \(E_{t}\) 包含两个不同的可比较元素。
\item
  每个包含至少两个不同的可比较元素的 \(E_{t}\) 是无间隙的。
\item
  如果 \(\alpha, \beta\) 是 \(I\) 中满足 \(\alpha < \beta\) 的两个元素，并且如果区间 \([]\alpha, \beta[\) （\(I\) 中的区间）为空，则要么 \(E_{\alpha}\) 没有极大元素，要么 \(E_{\beta}\) 没有极小元素。
\end{enumerate}

特别是，无间隙集合的每一个序数和 \(\sum_{i\in I}E_i\) 本身是无间隙的，只要其中没有 \(E_i\) 具有极大元素（或只要没有 \(E_i\) 具有极小元素）。如果 I 是无间隙的，并且每个 \(E_i\) 要么是无间隙的，要么不包含两个不同的可比较元素，则 \(\sum E_i\) 是无间隙的。

\begin{enumerate}
\def\labelenumi{\arabic{enumi}.}
\setcounter{enumi}{19}
\tightlist
\item
  一个有序集 E 被称为离散的（scattered），如果 E 的有序子集都不是无间隙的（练习 19）。离散集的每个子集都是离散的。* 每个多于一个元素的良序集都是离散的。*
\end{enumerate}

\begin{enumerate}
\def\labelenumi{(\alph{enumi})}
\item
  假设E是离散的。那么，如果x，y是E中任意两个元素，且x \textless{} y，存在两个元素 \(x'\) , \(y'\) 使得 \(x \leqslant x' < y' \leqslant y\) 并且使得区间 \(]x', y'[\) 是空的。* 给出一个满足此条件但不是离散的全序集的例子（考虑康托尔三分集）。*
\item
  序数和 \(E = \sum_{i \in I} E_i\) （其中既没有 I 也没有任何 \(E_i\) 为空）是离散的，当且仅当 I 和每个 \(E_i\) 都是离散的。（注意 E 包含一个与 I 同构的子集，并且 E 的每个子集 F 都是那些非空的 \(F \cap E_i\) 的序数和；最后使用习题 19。）
\end{enumerate}

\begin{enumerate}
\def\labelenumi{\arabic{enumi}.}
\setcounter{enumi}{20}
\item
  设E是一个非空全序集，并且令 \(S\{x,y\}\) 为关系“以x、y为端点的闭区间是离散的”（习题20）。证明S是一个等价关系，并且它在x和y上与E上的序关系是弱相容的（习题2），关于S的等价类是离散集，并且商有序集E/S（习题2）要么无间隙，要么由单个元素组成。由此推出E同构于一个离散集的序数和，其指标集要么无间隙，要么由单个元素组成。
\item
  \begin{enumerate}
  \def\labelenumii{(\alph{enumii})}
  \tightlist
  \item
    设E是一个有序集。E的子集U被称为开集，如果对于每一个 \(x \in U\) ，U 包含区间 \([x, \rightarrow[\) 。一个开集 U 被称为正则的，如果不存在开集 \(V \supset U\) （不同于 U），使得 U 在 V 中是共尾的。证明每个开集 U 恰好在一个正则开集 \(\tilde{U} (*)\) 中是共尾的。映射 \(U \to \tilde{U}\) 是递增的。若 U、V 是两个开集且满足 \(U \cap V = \varnothing\) ，那么也有 \(\tilde{U} \cap \tilde{V} = \varnothing\) 。
  \end{enumerate}
\end{enumerate}

\begin{enumerate}
\def\labelenumi{(\alph{enumi})}
\setcounter{enumi}{1}
\item
  证明 E 的正则开子集之集 R(E)（按包含序）是一个完全布尔格（习题 17）。要使 R(E) 恰由两个元素组成，当且仅当 E 非空且右有向。
\item
  如果 \(\mathbf{F}\) 是 \(\mathbf{E}\) 的共尾子集，证明映射 \(\mathbf{U} \to \mathbf{U} \cap \mathbf{F}\) 是从 \(\mathbf{R}(\mathbf{E})\) 到 \(\mathbf{R}(\mathbf{F})\) 上的一个同构。
\item
  如果 \(E_{1}\) , \(E_{2}\) 是两个有序集，则 \(E_{1} \times E_{2}\) 中的每个开集具有形式 \(U_{1} \times U_{2}\) ，其中 \(U_{i}\) 在 \(E_{i}\) 中是开的 (i=1, 2)。集合 \(\mathbf{R}(\mathbf{E}_{1} \times \mathbf{E}_{2})\) 同构于 \(\mathbf{R}(\mathbf{E}_{1}) \times \mathbf{R}(\mathbf{E}_{2})\) 。
\end{enumerate}

\begin{enumerate}
\def\labelenumi{\arabic{enumi}.}
\setcounter{enumi}{22}
\tightlist
\item
  设 E 是一个有序集，并设 \(\mathbf{R}_{0}(\mathbf{E}) = \mathbf{R}(\mathbf{E}) - \{\phi\}\) （练习 22）。对每个 \(x \in E\) ，让 \(r(x)\) 表示其中区间 \([x, \rightarrow[\) （这是一个开集）是共尾的唯一正则开集。如此定义的映射 r 称为 E 到 \(\mathbf{R}_{0}(\mathbf{E})\) 的典范映射。赋予 \(\mathbf{R}_{0}(\mathbf{E})\) 与包含关系相反的序关系。
\end{enumerate}

\begin{enumerate}
\def\labelenumi{(\alph{enumi})}
\item
  证明映射 \(r\) 是递增的，并且 \(r(\mathbf{E})\) 在 \(\mathbf{R}_0(\mathbf{E})\) 中是共尾的。
\item
  一个有序集 E 被称为反定向的，如果典范映射 \(r: \mathbf{E} \to \mathbb{R}_{0}(\mathbf{E})\) 是单射。为此，必要且充分的条件是以下两个条件成立：
\end{enumerate}

\begin{enumerate}
\def\labelenumi{(\Roman{enumi})}
\item
  如果 \(x\) 和 \(y\) 是 \(\mathbf{E}\) 中满足 \(x < y\) 的两个元素，则存在 \(z \in \mathbf{E}\) 使得 \(x < z\) ，并且使得区间 \([y, \rightarrow[\) 和 \([z, \rightarrow[\) 不相交。
\item
  如果 \(x\) 和 \(y\) 是 \(\mathbf{E}\) 的两个不可比较元素，则要么存在 \(x' \geqslant x\) 使得区间 \([x', \rightarrow[\) 和 \([y, \rightarrow[\) 不相交，要么存在 \(y' \geqslant y\) 使得区间 \([x, \rightarrow[\) 和 \([y', \rightarrow[\) 不相交。
\end{enumerate}

\begin{enumerate}
\def\labelenumi{(\alph{enumi})}
\setcounter{enumi}{2}
\tightlist
\item
  证明：对于每个有序集 E， \(R_{0}(E)\) 是反定向的，并且从 \(R_{0}(E)\) 到 \(R_{0}(R_{0}(E))\) 中的典范映射是双射（使用练习 22 (a)）。
\end{enumerate}

* 24. (a) 一个有序集 E 被称为（右侧）分支的，如果对于每个 \(x \in \mathbf{E}\) 存在 E 中的 \(y, z\) 使得 \(x \leqslant y, x \leqslant z\) ，并且

区间 \([y, \rightarrow[\) 和 \([z, \rightarrow[\) 不相交。一个没有极大元素的反定向集（练习 23）是分支的。

\begin{enumerate}
\def\labelenumi{(\alph{enumi})}
\setcounter{enumi}{1}
\tightlist
\item
  设 E 为 R 中形如
\end{enumerate}

\[
[ k. 2 ^ {- n}, (k + 1) 2 ^ {- n} ] (0 \leqslant k <   2 ^ {n}),
\]

的区间集合，按关系 \(\supset\) 排序。证明 E 是反定向的且没有极大元素。

\begin{enumerate}
\def\labelenumi{(\alph{enumi})}
\setcounter{enumi}{2}
\item
  给出一个分支集的例子，其中不存在反定向的共尾子集。（取 (b) 中定义的集合 E 与一个不含可数共尾子集的良序集的乘积，并使用练习 22。）
\item
  给出一个有序集 E 的例子，它不是反定向的，但有一个反定向的共尾子集（注意序数和 \(\sum_{\xi\in E} F_{\xi}\) 包含一个与 E 同构的共尾子集）。*
\end{enumerate}

